\chapter{Perluasan Lapangan dan Teori Galois}\label{ch:b3:galois}

Dapatkah setiap persamaan diselesaikan dengan \omterm{def:b3:galois:radical}{radikal}, sebagaimana disarankan rumus
kuadrat dan rumus kubik Cardano? Dapatkah sebuah sudut dibagi tiga
dengan penggaris dan jangka? Kedua pertanyaan itu, yang terbuka berabad-abad,
dijawab --- secara negatif --- oleh satu gagasan Évariste \omterm{def:b3:galois:galois}{Galois}:
kaitkan pada setiap polinomial sebuah \emph{grup} berhingga berisi simetri
akarnya, lalu bacalah jawabannya pada grup itu. Bab ini membangun
kamusnya: \omterm{def:b3:galois:extension}{perluasan lapangan} beserta \omterm{def:b3:galois:extension}{derajatnya}, \omterm{thm:b3:galois:splitting}{lapangan pemecah} dan
\omterm{def:b3:galois:closure}{penutup aljabar}, \omterm{thm:b3:galois:finitefields}{lapangan berhingga} (sebuah teori yang lengkap --- beserta
kesiklikan $\mathbb F_q^\times$ yang telah dijanjikan), separabilitas, lalu
\omterm{thm:b3:galois:fundamental}{korespondensi Galois} itu sendiri, dengan bukti penuh. Panennya:
kemustahilan konstruksi klasik, struktur
lapangan siklotomik, dan ketakterselesaian persamaan berderajat lima dengan
\omterm{def:b3:galois:radical}{radikal} --- kesederhanaan $A_5$ pada \cref{ch:b3:groups} yang akhirnya
mengenai sasarannya.

\section{Perluasan, derajat, dan kealjabaran}

\begin{definition}\label{def:b3:galois:extension}
Sebuah \emph{perluasan lapangan}\index{perluasan lapangan} $L/K$ adalah lapangan
$L$ yang memuat $K$ sebagai sublapangan; lalu $L$ menjadi ruang vektor-$K$,
dan \emph{derajat}\index{derajat perluasan} $[L:K]$ adalah
dimensinya. Perluasan itu disebut \emph{berhingga} jika $[L:K] <
\infty$. \emph{Karakteristik} sebuah lapangan adalah pembangun
$\geq 0$ dari kernel $\Z \to K$, $n \mapsto n\cdot 1$: nilainya
$0$ atau sebuah bilangan prima $p$; sejalan dengan itu $K$ memuat sublapangan terkecil
(\emph{lapangan prima}) yang isomorfik dengan $\Q$ atau dengan $\mathbb
F_p = \Z/p\Z$.
\end{definition}

\begin{theorem}[Hukum menara]\label{thm:b3:galois:tower}
Jika $K \subseteq L \subseteq M$, maka $[M:K] = [M:L]\,[L:K]$: jika
$(e_i)$ basis $L$ atas $K$ dan $(f_j)$ basis $M$
atas $L$, maka $(e_if_j)$ merupakan basis $M$ atas
$K$.\index{hukum menara}
\end{theorem}

\begin{proof}
Membangun: $x \in M$ ditulis $x = \sum_j \lambda_jf_j$
($\lambda_j \in L$), dan setiap $\lambda_j = \sum_i \mu_{ij}e_i$
($\mu_{ij} \in K$): sehingga $x = \sum_{i,j}\mu_{ij}e_if_j$. Bebas:
$\sum_{i,j}\mu_{ij}e_if_j = 0$ ditulis ulang menjadi $\sum_j (\sum_i
\mu_{ij}e_i)f_j = 0$; jumlah di dalamnya terletak di $L$, sehingga lenyap
(sebab $f_j$ bebas atas $L$); lalu semua $\mu_{ij} = 0$ (sebab $e_i$
bebas atas $K$).
\end{proof}

\begin{definition}\label{def:b3:galois:algebraic}
Misalkan $L/K$ dan $\alpha \in L$. Jika ada $P \in K[X]$ tak nol dengan
$P(\alpha) = 0$, maka $\alpha$ disebut \emph{aljabar}\index{unsur
aljabar} atas $K$; pembangun monik $\pi_\alpha$ dari \omterm{def:b3:rings:ideal}{ideal}
$\{P : P(\alpha) = 0\}$ di $K[X]$ disebut \emph{polinomial
minimalnya}\index{polinomial minimal (teori lapangan)}, yaitu
polinomial \omterm{def:b3:rings:divisibility}{tak tereduksi} (sebab $\pi = QR$ dengan $Q(\alpha) = 0$ memaksa $R$
konstan menurut keminimalan \omterm{def:b3:galois:extension}{derajatnya}). Jika tidak, $\alpha$ disebut
\emph{transenden}. Kita tulis $K(\alpha)$ untuk sublapangan terkecil
$L$ yang memuat $K$ dan $\alpha$, serta $K[\alpha]$ untuk
subring terkecilnya.
\end{definition}

\begin{theorem}\label{thm:b3:galois:simple}
Jika $\alpha$ bersifat \omterm{def:b3:galois:algebraic}{aljabar} atas $K$ dengan $d = \deg\pi_\alpha$,
maka
\[
K(\alpha) = K[\alpha] \cong K[X]/(\pi_\alpha),
\qquad [K(\alpha) : K] = d,
\]
dengan basis $1, \alpha, \dots, \alpha^{d-1}$. Sebaliknya, jika
$[L:K] < \infty$, maka setiap $\alpha \in L$ bersifat \omterm{def:b3:galois:algebraic}{aljabar} dengan \omterm{def:b3:galois:extension}{derajat}
$\deg \pi_\alpha$ yang membagi $[L:K]$.
\end{theorem}

\begin{proof}
Pemasukan nilai $K[X] \to L$, $P \mapsto P(\alpha)$, berpeta
$K[\alpha]$ dan berkernel $(\pi_\alpha)$; karena $\pi_\alpha$
\omterm{def:b3:rings:divisibility}{tak tereduksi}, $K[X]/(\pi_\alpha)$ merupakan \emph{lapangan}
(\cref{prop:b3:rings:primemaximal}: di dalam \omterm{def:b3:rings:pidufd}{DIU} $K[X]$, \omterm{def:b3:rings:divisibility}{unsur
tak tereduksi} membangun \omterm{def:b3:rings:primemaximal}{ideal maksimal}), sehingga $K[\alpha]$ menjadi lapangan
yang memuat $K$ dan $\alpha$: jadi ia sama dengan $K(\alpha)$. Kelas
$1, X, \dots, X^{d-1}$ membentuk basis kuosiennya (lewat pembagian
Euclid), dan dari sinilah basis dan \omterm{def:b3:galois:extension}{derajatnya}. Sebaliknya jika $[L:K]
= n < \infty$: maka $1, \alpha, \dots, \alpha^n$ bergantung, sehingga memberikan
polinomial penghapus; lalu $[K(\alpha):K] =
\deg\pi_\alpha$ membagi $n$ menurut \omterm{thm:b3:galois:tower}{hukum menara}.
\end{proof}

\begin{corollary}\label{cor:b3:galois:algclosed}
Jika $\alpha, \beta$ bersifat \omterm{def:b3:galois:algebraic}{aljabar} atas $K$, maka demikian pula $\alpha \pm
\beta$, $\alpha\beta$, $\alpha/\beta$ (dengan $\beta \ne 0$): jadi
unsur $L$ yang \omterm{def:b3:galois:algebraic}{aljabar} atas $K$ membentuk sublapangan $L$.
Lebih lanjut kealjabaran bersifat transitif: \omterm{def:b3:galois:algebraic}{aljabar} atas yang \omterm{def:b3:galois:algebraic}{aljabar} tetaplah
\omterm{def:b3:galois:algebraic}{aljabar}.
\end{corollary}

\begin{proof}
$K(\alpha, \beta) = (K(\alpha))(\beta)$ berhingga atas
$K(\alpha)$ (sebab $\beta$ \omterm{def:b3:galois:algebraic}{aljabar} atas $K \subseteq K(\alpha)$) dan
$K(\alpha)/K$ berhingga: menurut \omterm{thm:b3:galois:tower}{hukum menara} $[K(\alpha,\beta):K] <
\infty$, sehingga setiap unsur $K(\alpha, \beta)$ --- termasuk
keempat unsur yang disebut --- bersifat \omterm{def:b3:galois:algebraic}{aljabar} (\cref{thm:b3:galois:simple}).
Ketransitifannya: jika $\beta$ \omterm{def:b3:galois:algebraic}{aljabar} atas $L$ dan $L/K$ bersifat
\omterm{def:b3:galois:algebraic}{aljabar}, maka koefisien $c_0, \dots, c_{m-1}$ dari
$\pi_{\beta/L}$ membangun perluasan berhingga $F = K(c_0, \dots,
c_{m-1})$ atas $K$ (lewat \omterm{thm:b3:galois:tower}{hukum menara} berulang), dan $F(\beta)/F$
berhingga: jadi $[F(\beta):K] < \infty$, sehingga $\beta$ \omterm{def:b3:galois:algebraic}{aljabar} atas
$K$.
\end{proof}

\begin{example}\label{ex:b3:galois:degrees}
$[\Q(\sqrt2):\Q] = 2$, $[\Q(\sqrt[3]2):\Q] = 3$ (sebab $X^3 - 2$
\omterm{def:b3:rings:divisibility}{tak tereduksi} menurut \omterm{thm:b3:rings:criteria}{Eisenstein}), $[\Q(\zeta_p):\Q] = p - 1$ untuk
$\zeta_p = \eu^{2\iu\pi/p}$ (sebab $\Phi_p$ \omterm{def:b3:rings:divisibility}{tak tereduksi},
\cref{ex:b3:rings:eisenstein}). \omterm{thm:b3:galois:tower}{Hukum menara} sudah menjadi
senjata: $\sqrt[3]2 \notin \Q(\sqrt2)$, sebab $3 \nmid 2$.
\end{example}

\section{Lapangan pemecah; penutup aljabar}

\begin{theorem}[Lapangan pemecah]\label{thm:b3:galois:splitting}
Misalkan $P \in K[X]$ tak konstan. Terdapat sebuah \emph{lapangan
pemecah}\index{lapangan pemecah} bagi $P$ atas $K$: yaitu perluasan $L =
K(\alpha_1, \dots, \alpha_n)$ yang dibangun oleh akar $P$ dan yang di dalamnya
$P$ terpecah atas faktor linear. Lapangan itu tunggal sampai pada
isomorfisma-$K$, dan $[L:K] \leq (\deg P)!$.
\end{theorem}

\begin{proof}
\emph{Keberadaannya}, lewat induksi pada $\deg P$: ambil sebuah faktor
\omterm{def:b3:rings:divisibility}{tak tereduksi} $Q$ dari $P$; lapangan $K_1 = K[X]/(Q)$ memuat akar
$\alpha_1 = \bar X$ dari $Q$, dan karenanya dari $P$; tulis $P = (X -
\alpha_1)P_1$ atas $K_1$ lalu terapkan induksi pada $P_1$ atas $K_1$;
\omterm{def:b3:galois:extension}{derajatnya} berkalian menjadi paling banyak $n(n-1)\cdots = n!$.

\emph{Ketunggalannya} menyusul dari \emph{lema perluasan isomorfisma}
yang lebih kuat: misalkan $\sigma \colon K \to K'$ sebuah isomorfisma,
$P \in K[X]$, $P^\sigma$ polinomial yang koefisiennya terpetakan,
dan $L, L'$ \omterm{thm:b3:galois:splitting}{lapangan pemecah} $P, P^\sigma$; maka $\sigma$ dapat diperluas
menjadi isomorfisma $L \to L'$. Induksi pada $[L:K]$: jika $P$ terpecah
di $K$, maka $L = K$, dan $L' = K'$ (sebab $P^\sigma$ terpecah di $K'$,
sedangkan $L'$ dibangun oleh akarnya). Jika tidak, pilihlah sebuah akar
$\alpha \in L \setminus K$ dari faktor \omterm{def:b3:rings:divisibility}{tak tereduksi} $Q$ dari $P$
dengan $\deg Q \geq 2$; lalu $Q^\sigma$ merupakan faktor \omterm{def:b3:rings:divisibility}{tak tereduksi}
$P^\sigma$ yang berakar $\beta \in L'$; maka
\[
K(\alpha) \cong K[X]/(Q) \xrightarrow{\ \sigma\ }
K'[X]/(Q^\sigma) \cong K'(\beta)
\]
memperluas $\sigma$ dengan $\alpha \mapsto \beta$. Sekarang $L$ merupakan
\omterm{thm:b3:galois:splitting}{lapangan pemecah} $P$ atas $K(\alpha)$, dan $L'$ \omterm{thm:b3:galois:splitting}{lapangan pemecah} $P^\sigma$
atas $K'(\beta)$, dengan $[L : K(\alpha)] < [L:K]$: induksinya
memperluasnya lebih jauh menjadi $L \to L'$.
\end{proof}

\begin{definition}\label{def:b3:galois:closure}
Sebuah lapangan $\Omega$ disebut \emph{tertutup secara aljabar}\index{lapangan
tertutup secara aljabar} jika setiap polinomial tak konstan di $\Omega[X]$
berakar di $\Omega$ (sehingga terpecah). Sebuah \emph{penutup
aljabar}\index{penutup aljabar} bagi $K$ adalah perluasan
\omterm{def:b3:galois:algebraic}{aljabar} $\bar K/K$ dengan $\bar K$ tertutup secara \omterm{def:b3:galois:algebraic}{aljabar}.
\end{definition}

\begin{theorem}[Steinitz]\label{thm:b3:galois:steinitz}
Setiap lapangan $K$ mempunyai \omterm{def:b3:galois:closure}{penutup aljabar}, yang tunggal sampai pada
isomorfisma-$K$.
\end{theorem}

\begin{proof}
\emph{Keberadaan (konstruksi Artin).} Misalkan $R = K[(X_f)_f]$
ring polinomial dengan satu variabel $X_f$ untuk setiap
$f \in K[X]$ monik tak konstan, dan $I$ \omterm{def:b3:rings:ideal}{ideal} yang dibangun oleh semua
$f(X_f)$. \omterm{def:b3:rings:ideal}{Ideal} $I$ bersifat sejati: sebab relasi $1 = \sum_{i=1}^r g_i\,
f_i(X_{f_i})$ hanya melibatkan berhingga polinomial; di dalam
\omterm{thm:b3:galois:splitting}{lapangan pemecah} persekutuan $E$ bagi $f_1\cdots f_r$ pilihlah akar $\alpha_i$ dari
$f_i$ lalu masukkan $X_{f_i} \mapsto \alpha_i$ (variabel lainnya
$\mapsto 0$): diperoleh $1 = 0$, mustahil. Misalkan $\mathfrak m \supseteq I$
maksimal (\cref{thm:b3:rings:krull}; lewat Zorn) dan $K_1 = R/\mathfrak
m$: sebuah \omterm{def:b3:galois:extension}{perluasan lapangan} $K$ yang di dalamnya setiap $f \in
K[X]$ tak konstan berakar, yakni $\bar X_f$, dan yang bersifat \omterm{def:b3:galois:algebraic}{aljabar} atas
$K$ (sebab ia dibangun oleh $\bar X_f$ yang masing-masing \omterm{def:b3:galois:algebraic}{aljabar}).
Iterasikan: $K \subseteq K_1 \subseteq K_2 \subseteq \cdots$, dengan
$K_{n+1}$ berlaku terhadap $K_n$ seperti $K_1$ terhadap $K$, lalu misalkan $\Omega =
\bigcup_n K_n$, sebuah lapangan. Sembarang $g \in \Omega[X]$ tak konstan mempunyai
koefisien yang berhingga banyaknya di suatu $K_n$; faktor \omterm{def:b3:rings:divisibility}{tak tereduksi}
$g$ atas $K_n$ berakar di $K_{n+1} \subseteq
\Omega$: jadi $\Omega$ \omterm{def:b3:galois:closure}{tertutup secara aljabar}, dan \omterm{def:b3:galois:algebraic}{aljabar} atas
$K$ (sebab setiap $K_n$ demikian, lewat ketransitifan,
\cref{cor:b3:galois:algclosed}): jadi $\Omega$ merupakan \omterm{def:b3:galois:closure}{penutup
aljabar}.

\emph{Ketunggalan.} Misalkan $\Omega, \Omega'$ dua \omterm{def:b3:galois:closure}{penutup
aljabar}. Tinjau himpunan pasangan $(E, \tau)$ dengan $K
\subseteq E \subseteq \Omega$ dan $\tau\colon E \to \Omega'$ berupa
pembenaman-$K$, yang diurutkan menurut perluasan; himpunan itu tak kosong (memuat $(K,
\mathrm{id})$) dan induktif (gabungan sebuah rantai), sehingga lema Zorn memberikan
$(E_0, \tau_0)$ yang maksimal. Jika $E_0 \neq \Omega$, pilihlah $\alpha \in
\Omega\setminus E_0$: maka $\pi_{\alpha/E_0}$ terpetakan menjadi polinomial
atas $\tau_0(E_0)$ yang berakar $\beta$ di dalam $\Omega'$ yang \omterm{def:b3:galois:closure}{tertutup secara
aljabar}, sehingga $\tau_0$ dapat diperluas ke $E_0(\alpha) \to
\Omega'$ (dengan $\alpha \mapsto \beta$), dan itu bertentangan dengan kemaksimalan. Jadi
ada pembenaman-$K$ $\tau \colon \Omega \to \Omega'$; petanya,
yang isomorfik dengan $\Omega$, \omterm{def:b3:galois:closure}{tertutup secara aljabar}, dan
$\Omega'$ bersifat \omterm{def:b3:galois:algebraic}{aljabar} atasnya: sebab untuk $x \in \Omega'$,
$\pi_{x/\tau(\Omega)}$ terpecah atas $\tau(\Omega)$, sehingga $x \in
\tau(\Omega)$. Dengan demikian $\tau$ surjektif: jadi sebuah isomorfisma.
\end{proof}

\begin{remark}\label{rem:b3:galois:closureexamples}
Untuk $K = \Q$ mesin transfinit itu dapat dihindari:
bilangan \omterm{def:b3:galois:algebraic}{aljabar} $\bar\Q = \{z \in \C : z \text{ aljabar atas
} \Q\}$ membentuk sebuah \omterm{def:b3:galois:closure}{penutup aljabar} --- yaitu sublapangan $\C$ menurut
\cref{cor:b3:galois:algclosed}, yang \omterm{def:b3:galois:closure}{tertutup secara aljabar} karena $\C$
demikian (d'Alembert--Gauss, yang dibuktikan lewat analisis kompleks pada
\cref{ch:b3:holomorphic}) dan karena akar polinomial atas $\bar\Q$
bersifat \omterm{def:b3:galois:algebraic}{aljabar} atas $\Q$ menurut ketransitifan.
\end{remark}

\section{Lapangan berhingga}

\begin{theorem}\label{thm:b3:galois:finitefields}
Misalkan $p$ prima, $n \geq 1$, dan $q = p^n$.
\begin{enumerate}
  \item \omterm{thm:b3:galois:finitefields}{Lapangan berhingga} berkardinalitas pangkat prima, dan untuk
        setiap $q$ ada tepat satu lapangan $\mathbb F_q$ dengan
        $q$ unsur sampai isomorfisma: yaitu \omterm{thm:b3:galois:splitting}{lapangan pemecah}
        $X^q - X$ atas $\mathbb F_p$.\index{lapangan berhingga}
  \item \emph{Frobenius}\index{endomorfisma Frobenius} $F
        \colon x \mapsto x^p$ merupakan automorfisma $\mathbb
        F_q$, dan grup automorfisma $\mathbb F_q$ bersifat
        siklik berorde $n$, dibangun oleh $F$.
  \item $\mathbb F_{p^m}$ terbenam di $\mathbb F_{p^n}$ jika dan hanya jika $m \mid
        n$.
\end{enumerate}
\end{theorem}

\begin{proof}
(1) \omterm{thm:b3:galois:finitefields}{Lapangan berhingga} $E$ berkarakteristik $p > 0$ dan merupakan
ruang vektor-$\mathbb F_p$ berdimensi berhingga: jadi $\abs E = p^n$.
Grup multiplikatifnya berorde $q - 1$, sehingga setiap $x \in E$
memenuhi $x^q = x$: berarti $E$ terdiri atas $q$ akar $X^q - X$,
sehingga menjadi \omterm{thm:b3:galois:splitting}{lapangan pemecahnya} atas $\mathbb F_p$ ---
dan ini menentukan $E$ sampai isomorfisma
(\cref{thm:b3:galois:splitting}). Sebaliknya, di dalam \omterm{thm:b3:galois:splitting}{lapangan
pemecah} $L$ bagi $X^q - X$, himpunan $E$ berisi akarnya merupakan
\emph{sublapangan}: sebab $(x + y)^q = x^q + y^q$ lewat pengiterasian
mimpi mahasiswa baru $(a+b)^p = a^p + b^p$ (karena $p \mid \binom pk$), dan
$(xy)^q = x^qy^q$, $(x^{-1})^q = (x^q)^{-1}$; unsurnya tepat $q$
sebab $X^q - X$ bersifat \omterm{def:b3:galois:separable}{separabel}: turunannya
$qX^{q-1} - 1 = -1$ (sebab $p \mid q$), yang relatif prima dengannya, sehingga tak ada
akar berulang. Jadi $L = E$ beranggota $q$ unsur.

(2) $F$ adalah morfisma lapangan (lewat mimpi mahasiswa baru), yang injektif
(karena lapangan), sehingga bijektif pada $\mathbb F_q$ yang berhingga. Berlaku $F^n =
\mathrm{id}$ (sebab $x^q = x$), dan tak ada pangkat lebih kecil yang menjadi identitas:
sebab $F^m = \mathrm{id}$ berarti seluruh $q$ unsurnya menjadi akar
$X^{p^m} - X$, yang memaksa $p^m \geq q$. Jadi $\langle F\rangle$ bersifat
siklik berorde $n$; dan tak ada automorfisma lain, menurut
batas $\abs{\operatorname{Aut}} \leq [\,\mathbb F_q :
\mathbb F_p\,] = n$ yang dibuktikan di bawah
(\cref{prop:b3:galois:embeddings} dengan $L = \mathbb F_q$, $K =
\mathbb F_p$: sebab automorfisma menetapkan lapangan primanya).

(3) Jika $\mathbb F_{p^m} \subseteq \mathbb F_{p^n}$, maka \omterm{thm:b3:galois:tower}{hukum menara}
memberikan $p^n = (p^m)^d$: jadi $m \mid n$. Sebaliknya jika $m \mid n$, maka
$p^m - 1 \mid p^n - 1$ (lewat jumlah geometri), sehingga $X^{p^m} - X$ membagi
$X^{p^n} - X$ (argumen yang sama pada eksponennya: $X^{a} - 1 \mid X^{b}
- 1$ bila $a \mid b$), dan akar polinomial pertama yang berada di dalam
$\mathbb F_{p^n}$ membentuk sublapangan yang diminta, berkardinalitas
$p^m$ (dengan separabilitas seperti pada (1)).
\end{proof}

\begin{theorem}[Kesiklikan]\label{thm:b3:galois:cyclic}
Setiap subgrup berhingga dari grup multiplikatif sebuah lapangan bersifat
siklik. Khususnya $\mathbb F_q^\times \cong
\Z/(q-1)\Z$.\index{grup multiplikatif lapangan berhingga}
\end{theorem}

\begin{proof}
Misalkan $G \leq K^\times$ berhingga. Menurut teorema struktur
(\cref{cor:b3:modules:abelian}), $G \cong \Z/d_1 \times \dots
\times \Z/d_s$ dengan $d_1 \mid \dots \mid d_s$. Setiap $x \in G$
lalu memenuhi $x^{d_s} = 1$; padahal $X^{d_s} - 1$ berakar paling banyak
$d_s$ di lapangan $K$: sehingga $\abs G = d_1\cdots d_s \leq d_s$,
dan itu memaksa $s = 1$: jadi $G$ siklik.
\end{proof}

\begin{example}\label{ex:b3:galois:f8}
$\mathbb F_8 = \mathbb F_2[X]/(X^3 + X + 1)$: polinomial kubik itu tak
berakar di $\mathbb F_2$, sehingga \omterm{def:b3:rings:divisibility}{tak tereduksi}. Dengan menulis $\omega =
\bar X$: $\mathbb F_8^\times$ siklik berorde $7$, sehingga
\emph{setiap} unsur $\neq 0, 1$ membangunnya. Sublapangan
$\mathbb F_{p^{12}}$ membentuk kisi pembagi $12$:
$\mathbb F_p, \mathbb F_{p^2}, \mathbb F_{p^3}, \mathbb F_{p^4},
\mathbb F_{p^6}, \mathbb F_{p^{12}}$ --- contoh pertama yang lengkap
bagi \omterm{thm:b3:galois:fundamental}{korespondensi Galois}.
\end{example}

\section{Separabilitas dan pembenaman}

\begin{definition}\label{def:b3:galois:separable}
Sebuah polinomial $P \in K[X]$ disebut \emph{separabel}\index{polinomial
separabel} jika ia tak berakar berulang di sebuah \omterm{thm:b3:galois:splitting}{lapangan pemecah} ---
setara dengan $\gcd(P, P') = 1$ (sebab akar berulang menjadi akar
persekutuan; sebaliknya, atas \omterm{thm:b3:galois:splitting}{lapangan pemecahnya}, akar persekutuan bersifat
berulang; dan FPB-nya tak berubah di bawah \omterm{def:b3:galois:extension}{perluasan lapangan}, lewat
argumen \cref{cor:b3:modules:descent}). \omterm{def:b3:galois:algebraic}{Unsur aljabar} disebut
separabel jika \omterm{def:b3:galois:algebraic}{polinomial minimalnya} separabel; dan perluasan $L/K$ disebut
separabel jika semua unsurnya separabel.
\end{definition}

\begin{proposition}\label{prop:b3:galois:perfect}
Polinomial \emph{\omterm{def:b3:rings:divisibility}{tak tereduksi}} $P \in K[X]$ bersifat \omterm{def:b3:galois:separable}{separabel} kecuali jika $P' = 0$,
yang memaksa $\operatorname{char} K = p > 0$ dan $P \in K[X^p]$.
Akibatnya setiap perluasan \omterm{def:b3:galois:algebraic}{aljabar} dari lapangan berkarakteristik
$0$, dan dari \omterm{thm:b3:galois:finitefields}{lapangan berhingga}, bersifat \omterm{def:b3:galois:separable}{separabel} (lapangan
semacam itu disebut \emph{sempurna}\index{lapangan sempurna}).
\end{proposition}

\begin{proof}
FPB $\gcd(P, P')$ membagi $P$; jika nilainya bukan $1$, maka ketaktereduksian
memaksa $\gcd = P$ (sampai pada sebuah konstanta), sehingga $P \mid P'$ dengan $\deg
P' < \deg P$: jadi $P' = 0$. Dengan menulis $P = \sum a_kX^k$: berlaku $ka_k = 0$
untuk setiap $k$, sehingga pada karakteristik $0$, $P$ menjadi konstan (yang dikecualikan);
sedangkan pada karakteristik $p$, $a_k = 0$ kecuali bila $p \mid k$: jadi $P =
Q(X^p)$. Atas \omterm{thm:b3:galois:finitefields}{lapangan berhingga}, setiap unsur merupakan pangkat ke-$p$
(sebab \omterm{thm:b3:galois:finitefields}{Frobenius} surjektif), sehingga $Q(X^p) = \sum b_k^p X^{pk} = (\sum
b_kX^k)^p$ tidak \omterm{def:b3:rings:divisibility}{tak tereduksi}: jadi $P' = 0$ juga tak mungkin terjadi bagi
$P$ yang \omterm{def:b3:rings:divisibility}{tak tereduksi} di sana.
\end{proof}

\begin{proposition}[Mencacah pembenaman]\label{prop:b3:galois:embeddings}
Misalkan $L = K(\alpha_1, \dots, \alpha_r)$ berhingga atas $K$, dan
$\sigma \colon K \to \Omega$ sebuah pembenaman ke dalam lapangan yang \omterm{def:b3:galois:closure}{tertutup
secara aljabar}. Maka banyaknya perluasan $\sigma$ ke $L$
paling banyak $[L:K]$, dengan kesamaan bila $L/K$ \omterm{def:b3:galois:separable}{separabel}. Khususnya
$\abs{\operatorname{Aut}_K(L)} \leq [L:K]$.
\end{proposition}

\begin{proof}
Induksi pada $[L:K]$ lewat langkah sederhana. Untuk $L = K(\alpha)$: sebuah
perluasan $\tau$ ditentukan oleh $\tau(\alpha)$, yang haruslah sebuah
akar di $\Omega$ dari $\pi_\alpha^\sigma$; sebaliknya setiap akar demikian
memberikan satu perluasan (sebab $K(\alpha) \cong K[X]/(\pi_\alpha)$).
Banyaknya perluasan sama dengan banyaknya akar \emph{berbeda}
$\pi_\alpha^\sigma$ di $\Omega$: paling banyak $\deg\pi_\alpha =
[K(\alpha):K]$, dengan kesamaan jika dan hanya jika $\pi_\alpha$ \omterm{def:b3:galois:separable}{separabel}
(separabilitas $\pi^\sigma$ dan $\pi$ bersesuaian: FPB dengan
turunannya terpelihara oleh $\sigma$). Secara umum, faktorkan $L =
K(\alpha_1)(\alpha_2, \dots)$: perluasan $\sigma$ ke
$K(\alpha_1)$ berjumlah $\leq [K(\alpha_1):K]$, dan masing-masing meluas dalam
$\leq [L : K(\alpha_1)]$ cara menurut induksi; lalu kalikan (lewat \omterm{thm:b3:galois:tower}{hukum
menara}). Pada kasus \omterm{def:b3:galois:separable}{separabel} kedua cacahan itu menjadi kesamaan: sebab polinomial
minimal atas lapangan $K(\alpha_1)$ yang lebih besar membagi polinomial minimal
atas $K$, sehingga tetap \omterm{def:b3:galois:separable}{separabel}.
\end{proof}

\begin{theorem}[Unsur primitif]\label{thm:b3:galois:primitive}
Setiap perluasan \emph{\omterm{def:b3:galois:separable}{separabel}} yang berhingga bersifat sederhana: $L =
K(\gamma)$ untuk suatu $\gamma$.\index{teorema unsur primitif}
\end{theorem}

\begin{proof}
Jika $K$ berhingga, maka $L$ pun berhingga, dan sebuah pembangun $\gamma$ dari
grup siklik $L^\times$ (\cref{thm:b3:galois:cyclic}) sudah mencukupi.
Misalkan $K$ tak berhingga; lewat induksi cukup ditangani kasus $L =
K(\alpha, \beta)$. Misalkan $n = [L:K]$; menurut
\cref{prop:b3:galois:embeddings} ada $n$ pembenaman-$K$ berbeda
$\sigma_1, \dots, \sigma_n \colon L \to \Omega$
(dengan $\Omega$ sebuah \omterm{def:b3:galois:closure}{penutup aljabar}). Polinomial
\[
D(T)=\prod_{i<j}\bigl[\bigl(\sigma_i(\alpha)-\sigma_j(\alpha)
\bigr) + T\bigl(\sigma_i(\beta) - \sigma_j(\beta)\bigr)\bigr]
\]
tidak nol secara identik: sebab sebuah faktornya lenyap secara identik hanya jika
$\sigma_i, \sigma_j$ bersesuaian baik pada $\alpha$ maupun $\beta$, sehingga pada
$L$ --- dan itu dikecualikan untuk $i \neq j$. Karena $K$ tak berhingga, pilihlah $c \in
K$ dengan $D(c) \ne 0$: maka $n$ unsur $\sigma_i(\alpha +
c\beta)$ berbeda berpasangan, sehingga $\gamma = \alpha + c\beta$
mempunyai paling sedikit $n$ konjugat berbeda di $\Omega$, yakni
$\deg \pi_\gamma \geq n$: dan $[K(\gamma):K] \geq n = [L:K]$ memaksa
$L = K(\gamma)$.
\end{proof}

\section{Korespondensi Galois}

\begin{definition}\label{def:b3:galois:galois}
Perluasan berhingga $L/K$ disebut \emph{Galois}\index{perluasan Galois}
jika ia merupakan \omterm{thm:b3:galois:splitting}{lapangan pemecah} sebuah polinomial \emph{\omterm{def:b3:galois:separable}{separabel}}
atas $K$. \emph{Grup Galois}\index{grup Galois} $L/K$ adalah
$\operatorname{Gal}(L/K) = \operatorname{Aut}_K(L)$, yaitu grup
automorfisma lapangan $L$ yang menetapkan $K$ titik demi titik.
\end{definition}

\begin{proposition}\label{prop:b3:galois:galoisorder}
Jika $L/K$ bersifat \omterm{def:b3:galois:galois}{Galois}, maka $\abs{\operatorname{Gal}(L/K)} =
[L:K]$; lebih lanjut $L/F$ bersifat \omterm{def:b3:galois:galois}{Galois} untuk setiap lapangan antara $K
\subseteq F \subseteq L$, dan setiap pembenaman-$F$ $L \to \Omega
\supseteq L$ berpeta $L$ (\emph{kenormalan}).
\end{proposition}

\begin{proof}
Misalkan $L$ memecah \omterm{def:b3:galois:separable}{polinomial separabel} $P$ atas $K$, dan tetapkan sebuah
\omterm{def:b3:galois:closure}{penutup aljabar} $\Omega \supseteq L$. Perluasan $L/K$ bersifat \omterm{def:b3:galois:separable}{separabel}: sebab ia dibangun
oleh akar $P$; separabilitas setiap unsurnya menyusul dari
kasus kesamaan di bawah, tetapi mari kita berargumen langsung ---
\cref{prop:b3:galois:embeddings} yang diterapkan pada pembangunnya (yaitu akar
\omterm{def:b3:galois:separable}{polinomial separabel} $P$, yang \omterm{def:b3:galois:algebraic}{polinomial minimalnya} membagi $P$)
melahirkan tepat $[L:K]$ perluasan dari $K \hookrightarrow \Omega$
(pada langkah induktifnya, polinomial minimal sebuah akar $P$
atas lapangan antara tetap membagi $P$, sehingga tetap
\omterm{def:b3:galois:separable}{separabel}). Setiap pembenaman $\tau \colon L \to \Omega$ semacam itu
mempermutasikan akar $P$ (sebab $\tau$ menetapkan koefisiennya), dan
$L$ dibangun olehnya: jadi $\tau(L) = L$. Karena itu pembenaman $=$
automorfisma: $\abs{\operatorname{Gal}(L/K)} = [L:K]$. Untuk
lapangan antara $F$: $L$ juga \omterm{thm:b3:galois:splitting}{lapangan pemecah} $P$ atas
$F$, dan $P$ tetap \omterm{def:b3:galois:separable}{separabel}: jadi $L/F$ bersifat \omterm{def:b3:galois:galois}{Galois}; argumen yang sama
memberikan kenormalan atas $F$.
\end{proof}

\begin{lemma}[Artin]\label{lem:b3:galois:artin}
Misalkan $G$ grup berhingga berisi automorfisma sebuah lapangan $L$ dan
$K = L^G = \{x : \sigma(x) = x\ \forall\sigma \in G\}$ lapangan
tetapnya. Maka $[L : L^G] \leq \abs G$.
\end{lemma}

\begin{proof}
Misalkan $n = \abs G$, $G = \{\sigma_1, \dots, \sigma_n\}$, dan
andaikan $x_1, \dots, x_{n+1} \in L$ bebas linear atas
$K$. Sistem linear homogen dengan $n$ persamaan dan $n+1$
variabel $(c_j)$ atas $L$,
\[
\sum_{j=1}^{n+1} c_j\,\sigma_i(x_j) = 0
\qquad (i = 1, \dots, n),
\]
mempunyai penyelesaian tak nol; pilihlah satu yang entri tak nolnya
\emph{paling sedikit}, katakanlah $c_1, \dots, c_r \neq 0$ (setelah dinomori ulang), $r \geq 2$
(sebab satu $c_j\sigma_i(x_j) = 0$ tunggal mustahil), yang dinormalkan
$c_r = 1$. Tidak semua $c_j$ terletak di $K$: sebab persamaan untuk $\sigma_i =
\mathrm{id}$ akan bertentangan dengan kebebasannya; katakanlah $c_1 \notin K$,
sehingga $\tau(c_1) \ne c_1$ untuk suatu $\tau \in G$. Terapkan $\tau$ pada semua
persamaannya: karena $\tau\sigma_i$ menjelajahi $G$, vektor
$(\tau(c_j))_j$ merupakan penyelesaian lain; setelah dikurangkan, $(c_j -
\tau(c_j))_j$ menjadi penyelesaian dengan entri tak nol yang lebih sedikit (sebab
entri ke-$r$ menjadi $1 - 1 = 0$, sedangkan entri pertamanya tidak) dan tidak
nol: kontradiksi. Jadi sembarang $n+1$ unsur bergantung: $[L:K]
\leq n$.
\end{proof}

\begin{theorem}[Teorema fundamental teori Galois]
\label{thm:b3:galois:fundamental}
Misalkan $L/K$ sebuah \omterm{def:b3:galois:galois}{perluasan Galois} dengan grup $G =
\operatorname{Gal}(L/K)$.\index{korespondensi Galois}
\begin{enumerate}
  \item $L^G = K$.
  \item Pemetaan $H \mapsto L^H$ dan $F \mapsto
        \operatorname{Gal}(L/F)$ saling invers, dan keduanya
        bijeksi yang membalik inklusi antara subgrup $G$
        dan lapangan antara $K \subseteq F \subseteq L$;
        lebih lanjut $[L : L^H] = \abs H$ dan $[L^H : K] = [G : H]$.
  \item $H \trianglelefteq G$ jika dan hanya jika $L^H/K$ bersifat \omterm{def:b3:galois:galois}{Galois}, dan lalu
        pembatasannya menginduksikan $\operatorname{Gal}(L^H/K) \cong
        G/H$.
\end{enumerate}
\end{theorem}

\begin{proof}
(1) Jelas $K \subseteq L^G$. Sebaliknya misalkan $\alpha \in L
\setminus K$; kita tampilkan $\sigma \in G$ dengan $\sigma(\alpha) \ne
\alpha$. Polinomial minimal $\pi_\alpha$ atas $K$ berderajat
$\geq 2$ dan bersifat \omterm{def:b3:galois:separable}{separabel} (sebab $L/K$ \omterm{def:b3:galois:separable}{separabel},
\cref{prop:b3:galois:galoisorder}), sehingga ia berakar lain
$\beta \neq \alpha$ di sebuah \omterm{def:b3:galois:closure}{penutup aljabar} $\Omega \supseteq L$.
Perluaslah pembenaman-$K$ $K(\alpha) \to \Omega$, $\alpha \mapsto
\beta$, menjadi pembenaman $\tau\colon L \to \Omega$
(\cref{prop:b3:galois:embeddings}); menurut kenormalan
(\cref{prop:b3:galois:galoisorder}) $\tau(L) = L$, sehingga $\tau \in
G$, $\beta = \tau(\alpha) \in L$, dan $\tau(\alpha) \neq
\alpha$.

(2) Untuk sebuah subgrup $H$: $L/L^H$ bersifat \omterm{def:b3:galois:galois}{Galois}
(\cref{prop:b3:galois:galoisorder}), dan
$\operatorname{Gal}(L/L^H) \supseteq H$ secara trivial, sehingga $[L:L^H] =
\abs{\operatorname{Gal}(L/L^H)} \geq \abs H$; sedangkan lema Artin
memberikan $[L:L^H] \leq \abs H$: jadi terjadi kesamaan, dan
$\operatorname{Gal}(L/L^H) = H$. Untuk lapangan antara $F$:
$L/F$ yang \omterm{def:b3:galois:galois}{Galois} memberikan $L^{\operatorname{Gal}(L/F)} = F$ menurut (1) yang diterapkan pada $L/F$.
Kedua pemetaan itu saling invers; dan jelas keduanya
membalik inklusi. \omterm{def:b3:galois:extension}{Derajatnya}: $[L:L^H] = \abs H$ baru saja
dibuktikan, dan $[L^H:K] = [L:K]/[L:L^H] = \abs G/\abs H$.

(3) Untuk $\sigma \in G$ dan $H \leq G$: $\sigma(L^H) =
L^{\sigma H\sigma^{-1}}$ (lewat pemeriksaan langsung). Menurut bijeksinya,
$\sigma(L^H) = L^H$ untuk setiap $\sigma$ jika dan hanya jika $H \trianglelefteq G$.
Sekarang jika $H \trianglelefteq G$, tetapkan $F = L^H$: setiap $\sigma \in
G$ terbatas menjadi automorfisma $F$, sehingga memberikan morfisma $\rho
\colon G \to \operatorname{Aut}_K(F)$ dengan kernel $\{\sigma :
\sigma\restriction_F = \mathrm{id}\} = \operatorname{Gal}(L/F) =
H$. Jadi $G/H$ terbenam di $\operatorname{Aut}_K(F)$, sehingga
$\abs{\operatorname{Aut}_K(F)} \geq [G:H] = [F:K]$; sedangkan ketaksamaan
sebaliknya selalu berlaku (\cref{prop:b3:galois:embeddings}):
jadi $\abs{\operatorname{Aut}_K(F)} = [F:K]$ dan $\rho$ surjektif. Tersisa
menunjukkan bahwa $F/K$ bersifat \omterm{def:b3:galois:galois}{Galois}: $F$ \omterm{def:b3:galois:separable}{separabel} atas $K$
(karena berada di dalam $L/K$ yang \omterm{def:b3:galois:separable}{separabel}), dan $F = K(\gamma)$
(\cref{thm:b3:galois:primitive}); polinomial
$\prod_{\sigma \in G/H}\bigl(X - \sigma(\gamma)\bigr)$ (hasil kali
atas peta yang berbeda, yang terletak di $F$ sebab $\sigma(F) =
L^{\sigma H\sigma^{-1}} = L^H = F$ menurut kenormalan $H$) mempunyai
koefisien
yang ditetapkan $G$, sehingga terletak di $K$ menurut (1): jadi ia \omterm{def:b3:galois:separable}{polinomial separabel}
di $K[X]$ yang dipecah oleh $F$, dan akarnya membangun $F$: jadi $F/K$ bersifat
\omterm{def:b3:galois:galois}{Galois}. Sebaliknya, jika $F = L^H$ dengan $F/K$ \omterm{def:b3:galois:galois}{Galois}, maka kenormalan
$F$ (\cref{prop:b3:galois:galoisorder}, yang diterapkan pada pembenaman
$F \to \Omega$ hasil pembatasan unsur $G$) memberikan
$\sigma(F) = F$ untuk setiap $\sigma \in G$, yakni $H
\trianglelefteq G$.
\end{proof}

\begin{omfigure}
\begin{tikzpicture}[xscale=1.5, yscale=1.15]
  \tikzset{fd/.style={font=\small}}
  % subgroup lattice of S3 (left, upside down to mirror fields)
  \node[fd] (G)  at (-3.2, 0)   {$S_3$};
  \node[fd] (A)  at (-4.4, 1.2) {$\langle(1\,2\,3)\rangle$};
  \node[fd] (T1) at (-3.6, 1.2) {$\langle(1\,2)\rangle$};
  \node[fd] (T2) at (-2.8, 1.2) {$\langle(1\,3)\rangle$};
  \node[fd] (T3) at (-2.0, 1.2) {$\langle(2\,3)\rangle$};
  \node[fd] (E)  at (-3.2, 2.4) {$\{e\}$};
  \draw (G) -- (A); \draw (G) -- (T1); \draw (G) -- (T2);
  \draw (G) -- (T3);
  \draw (E) -- (A); \draw (E) -- (T1); \draw (E) -- (T2);
  \draw (E) -- (T3);
  % field lattice (right)
  \node[fd] (Q)  at (1.6, 0)   {$\Q$};
  \node[fd] (W)  at (0.4, 1.2) {$\Q(\iu\sqrt3)$};
  \node[fd] (R1) at (1.4, 1.2) {$\Q(\sqrt[3]2)$};
  \node[fd] (R2) at (2.4, 1.2) {$\Q(j\sqrt[3]2)$};
  \node[fd] (R3) at (3.5, 1.2) {$\Q(j^2\sqrt[3]2)$};
  \node[fd] (L)  at (1.6, 2.4) {$L = \Q(\sqrt[3]2, j)$};
  \draw (Q) -- (W); \draw (Q) -- (R1); \draw (Q) -- (R2);
  \draw (Q) -- (R3);
  \draw (L) -- (W); \draw (L) -- (R1); \draw (L) -- (R2);
  \draw (L) -- (R3);
  \draw[->, omDef, thick, shorten >=6pt, shorten <=6pt]
    (-1.4, 1.2) -- (-0.6, 1.2);
\end{tikzpicture}

{\small \omterm{thm:b3:galois:fundamental}{Korespondensi Galois} untuk \omterm{thm:b3:galois:splitting}{lapangan pemecah} $L$ dari
$X^3 - 2$ atas $\Q$ ($j = \eu^{2\iu\pi/3}$): subgrup
$\operatorname{Gal}(L/\Q) \cong S_3$ (kiri, urutannya dibalik)
berpadanan dengan lapangan antara (kanan). Satu-satunya \omterm{def:b3:groups:normal}{subgrup normal}
sejati $\langle(1\,2\,3)\rangle$ bersesuaian dengan satu-satunya
subperluasan $\Q(\iu\sqrt3)/\Q$ yang bersifat \omterm{def:b3:galois:galois}{Galois}; sedangkan ketiga
subgrup sekawan $\langle(i\,j)\rangle$ bersesuaian dengan ketiga
lapangan kubik sekawan $\Q(j^k\sqrt[3]2)$, yang tak satu pun
normal atas $\Q$.}
\end{omfigure}

\section{Perluasan siklotomik}

\begin{definition}\label{def:b3:galois:cyclotomic}
Misalkan $n \geq 1$ dan $\zeta_n = \eu^{2\iu\pi/n}$.
Yang disebut \emph{polinomial siklotomik}\index{polinomial siklotomik} ke-$n$ adalah
$\Phi_n = \prod_{\gcd(k,n)=1,\ 1 \le k \le n} \bigl(X -
\zeta_n^k\bigr)$, yang berderajat $\varphi(n)$; dengan mengelompokkan akar
$X^n - 1$ menurut orde eksaknya, $X^n - 1 = \prod_{d \mid n}\Phi_d$,
dan ini menunjukkan secara induktif bahwa $\Phi_n \in \Z[X]$ (lewat pembagian
Euclid polinomial bulat yang monik).
\end{definition}

\begin{theorem}\label{thm:b3:galois:cyclotomicirred}
Polinomial $\Phi_n$ \omterm{def:b3:rings:divisibility}{tak tereduksi} atas $\Q$; karena itu $[\Q(\zeta_n) : \Q] =
\varphi(n)$ dan
\[
\operatorname{Gal}\bigl(\Q(\zeta_n)/\Q\bigr) \;\cong\;
(\Z/n\Z)^\times,
\qquad \sigma_a(\zeta_n) = \zeta_n^a .
\]
Dengan demikian perluasan $\Q(\zeta_n)/\Q$ bersifat \omterm{def:b3:galois:galois}{Galois} dengan grup
\emph{abelian}.
\end{theorem}

\begin{proof}
Misalkan $f = \pi_{\zeta_n}$, sehingga $\Phi_n = fg$ dengan $f, g \in \Z[X]$
monik (lewat lema Gauss \cref{lem:b3:rings:gauss}: kandungannya
berkalian, dan semua polinomialnya monik). \emph{Klaim: jika $\zeta$
akar $f$ dan $p \nmid n$ bilangan prima, maka $\zeta^p$ juga akar
$f$.} Kalau tidak, $\zeta^p$ menjadi akar $g$ (sebab ia akar satuan
primitif ke-$n$), sehingga $\zeta$ menjadi akar $g(X^p)$, dan $f
\mid g(X^p)$ di $\Z[X]$ (lewat polinomial minimal, lalu lema Gauss lagi).
Reduksikan modulo $p$: $\bar g(X^p) = \bar g(X)^p$ (lewat \omterm{thm:b3:galois:finitefields}{Frobenius} pada
$\mathbb F_p[X]$: koefisien demi koefisien $a^p = a$, ditambah mimpi
mahasiswa baru), sehingga $\bar f \mid \bar g^{\,p}$: jadi $\bar f$ dan $\bar g$
berbagi faktor \omterm{def:b3:rings:divisibility}{tak tereduksi}, dan $\bar\Phi_n = \bar f\bar g$ mempunyai
faktor berulang. Maka demikian pula $X^n - \bar 1$; padahal
turunannya $\bar nX^{n-1}$ relatif prima dengannya (sebab $p \nmid n$, dan
$0$ bukan akarnya): kontradiksi.

Setiap akar primitif $\zeta_n^k$ (dengan $\gcd(k, n) = 1$) diperoleh
dari $\zeta_n$ lewat pangkat prima berturutan yang tak membagi $n$
(faktorkan $k$): jadi klaimnya merambat, sehingga setiap akar primitif menjadi
akar $f$: berarti $f = \Phi_n$, yang \omterm{def:b3:rings:divisibility}{tak tereduksi}. Akibatnya
$[\Q(\zeta_n):\Q] = \varphi(n)$, dan $\Q(\zeta_n)$ merupakan
\omterm{thm:b3:galois:splitting}{lapangan pemecah} \omterm{def:b3:galois:separable}{polinomial separabel} $X^n - 1$ (sebab semua akarnya pangkat
$\zeta_n$): jadi \omterm{def:b3:galois:galois}{Galois}. Sebuah automorfisma $\sigma$ mengirim $\zeta_n$
ke akar primitif lain $\zeta_n^{a(\sigma)}$, dan $\sigma
\mapsto a(\sigma)$ merupakan morfisma injektif ke dalam
$(\Z/n\Z)^\times$; kedua grup itu berorde $\varphi(n)$: jadi
isomorfisma.
\end{proof}

\section{Penggaris dan jangka}

\begin{definition}\label{def:b3:galois:constructible}
Kenalilah bidang dengan $\C$; mulailah dari $\{0, 1\}$. Sebuah titik disebut
\emph{dapat dilukis}\index{bilangan yang dapat dilukis} jika ia
diperoleh lewat berhingga perpotongan garis melalui dua
titik yang sudah terlukis dan lingkaran berpusat di titik terlukis
dengan jari-jari berupa jarak dua titik terlukis.
\end{definition}

\begin{theorem}[Wantzel]\label{thm:b3:galois:wantzel}
Bilangan $z \in \C$ \omterm{def:b3:galois:constructible}{dapat dilukis} jika dan hanya jika ada menara $\Q = F_0
\subseteq F_1 \subseteq \dots \subseteq F_r$ dengan $[F_{i+1} :
F_i] = 2$ dan $z \in F_r$. Khususnya, \omterm{def:b3:galois:constructible}{bilangan yang dapat dilukis}
bersifat \omterm{def:b3:galois:algebraic}{aljabar} dengan \omterm{def:b3:galois:extension}{derajat} berupa pangkat $2$ atas
$\Q$.\index{teorema Wantzel}
\end{theorem}

\begin{proof}
($\Rightarrow$) Koordinat perpotongan dua garis
melalui titik yang koordinatnya di sublapangan $F \subseteq \R$
menyelesaikan sistem linear atas $F$: jadi koordinatnya tetap di $F$. Perpotongan garis--lingkaran
dan lingkaran--lingkaran menghasilkan, setelah bagian linearnya dilenyapkan
(mengurangkan kedua persamaan lingkaran memberikan sebuah garis),
sebuah persamaan kuadrat atas $F$: sehingga koordinat barunya terletak di $F$
atau di $F(\sqrt d)$ untuk suatu $d \in F$, $d > 0$. Lewat induksi,
setiap titik terlukis berkoordinat di sebuah menara perluasan
kuadrat atas $\Q$; dan $z = x + \iu y$ juga berada di menara kuadrat
(tambahkan $\iu$: satu langkah kuadrat lagi). Akibat pada \omterm{def:b3:galois:extension}{derajatnya}:
$[\Q(z):\Q]$ membagi $[F_r : \Q] = 2^r$ (lewat \omterm{thm:b3:galois:tower}{hukum
menara}).

($\Leftarrow$) \omterm{def:b3:galois:constructible}{Bilangan yang dapat dilukis} membentuk sebuah lapangan: jumlah dan
selisihnya lewat jajargenjang (garis sejajar \omterm{def:b3:galois:constructible}{dapat dilukis}:
turunkan lalu naikkan garis tegak lurus dua kali --- garis tegak lurus klasik
melalui sebuah titik memakai satu lingkaran dan dua busur);
hasil kali dan hasil baginya lewat konfigurasi perbandingan Thales
(diberikan panjang $a, b$, lukislah $ab$ dan $a/b$ dengan segitiga
sebangun pada dua sinar garis). Dan lapangan itu tertutup terhadap akar
kuadrat: untuk $a > 0$, lingkaran berdiameter $1 + a$ beserta
garis tegak lurus pada titik sambungnya berpotongan pada ketinggian $\sqrt a$
(hubungan rata-rata geometri pada garis tinggi segitiga siku-siku); sedangkan untuk
bilangan kompleks $w = \rho\eu^{\iu\theta}$, lukislah $\sqrt\rho$ lalu
bagilah dua $\theta$ (sudut dapat dibagi dua dengan jangka).
Karena itu bagian real dan imajiner unsur sebuah menara kuadrat
\omterm{def:b3:galois:constructible}{dapat dilukis} lewat induksi pada menaranya: sebab setiap langkahnya
menambahkan akar sebuah polinomial kuadrat, yang dapat dinyatakan lewat operasi lapangan
dan satu akar kuadrat dari bilangan yang sudah terlukis (lewat
rumus kuadrat; pada karakteristik $0$).
\end{proof}

\begin{corollary}\label{cor:b3:galois:impossible}
Ketiga soal klasik berikut tak \omterm{def:b3:groups:derived}{terselesaikan} dengan penggaris dan
jangka:\index{penggaris dan jangka}
\begin{enumerate}
  \item \emph{Penggandaan kubus}: $\sqrt[3]2$ berderajat
        $3$, bukan pangkat $2$.
  \item \emph{Pembagian tiga sudut}: membagi tiga $60^\circ$
        memerlukan $\cos 20^\circ$, yaitu akar polinomial \omterm{def:b3:rings:divisibility}{tak tereduksi} $8X^3
        - 6X - 1$: berderajat $3$.
  \item \emph{Pengkuadratan lingkaran}: $\sqrt\pi$ bersifat transenden
        (sebab $\pi$ demikian --- teorema Lindemann, yang di sini diterima tanpa bukti:
        pembuktiannya termasuk kuliah teori transendensi).
\end{enumerate}
Selain itu, segi-$n$ beraturan \omterm{def:b3:galois:constructible}{dapat dilukis} jika dan hanya jika
$\varphi(n)$ berupa pangkat $2$ (Gauss--Wantzel; arah ``jika'' memakai
metode soal akhir pekan, sedangkan arah ``hanya jika'' adalah
\cref{thm:b3:galois:wantzel} yang diterapkan pada $\zeta_n$, yang berderajat
$\varphi(n)$). Untuk $n = 7$: $\varphi(7) = 6$, sehingga segi tujuh
beraturan mustahil; untuk $n = 17$: $\varphi(17) = 16 = 2^4$, sehingga
\omterm{def:b3:galois:constructible}{dapat dilukis} --- dan soal akhir pekan melukisnya.
\end{corollary}

\begin{proof}
(1) $X^3 - 2$ \omterm{def:b3:rings:divisibility}{tak tereduksi} (menurut \omterm{thm:b3:rings:criteria}{Eisenstein}). (2) Dari $\cos
3\theta = 4\cos^3\theta - 3\cos\theta$ dengan $3\theta =
60^\circ$: diperoleh $8c^3 - 6c = 1$ untuk $c = \cos 20^\circ$; polinomial kubik
$8X^3 - 6X - 1$ tak berakar rasional (calonnya $\pm1, \pm\frac
12, \pm\frac14, \pm\frac18$ semuanya gagal), sehingga \omterm{def:b3:rings:divisibility}{tak tereduksi}: berderajat
$3$. Sudut $60^\circ$ umum \omterm{def:b3:galois:constructible}{dapat dilukis}, sehingga alat pembagi tiga
akan melukis $c$. (3) Seandainya $\sqrt\pi$ \omterm{def:b3:galois:constructible}{dapat dilukis}, ia
akan bersifat \omterm{def:b3:galois:algebraic}{aljabar}, sehingga $\pi$ pun demikian. Untuk pernyataan segi-$n$:
\omterm{def:b3:galois:extension}{derajat} $\zeta_n$ adalah $\varphi(n)$
(\cref{thm:b3:galois:cyclotomicirred}); keperluannya menyusul dari
Wantzel; sedangkan untuk kecukupannya, grup Galoisnya, yang abelian berorde
$2^m$, mempunyai rantai subgrup berindeks $2$ (dan grup-$2$ berhingga
memang demikian: \cref{exo:b3:groups:10}), yang lapangan tetapnya membentuk
menara kuadrat yang berakhir di $\Q(\zeta_n)$
(\cref{thm:b3:galois:fundamental}); lalu simpulkan dengan
\cref{thm:b3:galois:wantzel}.
\end{proof}

\section{Keterselesaian dengan radikal}

\begin{definition}\label{def:b3:galois:radical}
Sebuah perluasan $L/K$ (dengan karakteristik $0$ sepanjang bagian ini)
disebut \emph{radikal} jika ada menara $K = F_0 \subseteq \dots
\subseteq F_r = L$ dengan $F_{i+1} = F_i(\alpha_i)$ dan
$\alpha_i^{n_i} \in F_i$: jadi setiap langkahnya menambahkan akar ke-$n_i$. Sebuah
polinomial $P \in K[X]$ disebut \emph{terselesaikan dengan
radikal}\index{terselesaikan dengan radikal} jika \omterm{thm:b3:galois:splitting}{lapangan pemecahnya}
termuat di suatu perluasan radikal atas $K$.
\end{definition}

\begin{lemma}\label{lem:b3:galois:cyclicsteps}
Misalkan $K$ memuat akar satuan primitif ke-$n$ bernama $\zeta$, yakni
$\zeta$ berorde $n$ di $K^\times$, dan $a \in K^\times$. Maka
$K(\sqrt[n]a)/K$ bersifat \omterm{def:b3:galois:galois}{Galois} dengan grup \emph{siklik}. Sebaliknya
--- meski tak diperlukan di bawah --- setiap perluasan siklik berderajat $n$
berbentuk demikian. Lebih lanjut $K(\zeta_n)/K$ bersifat \omterm{def:b3:galois:galois}{Galois} dengan grup
\emph{abelian}, untuk $K$ berkarakteristik $0$ mana pun.
\end{lemma}

\begin{proof}
Polinomial $X^n - a$ bersifat \omterm{def:b3:galois:separable}{separabel} (lihat $\gcd$ dengan $nX^{n-1}$: sebab $a \neq 0$) dan
terpecah di $K(\alpha)$ dengan $\alpha^n = a$: akarnya adalah
$\zeta^k\alpha \in K(\alpha)$. Jadi $K(\alpha)/K$ bersifat \omterm{def:b3:galois:galois}{Galois};
pemetaan $\sigma \mapsto \sigma(\alpha)/\alpha \in \mu_n = \langle
\zeta\rangle$ merupakan morfisma injektif (sebab $\sigma\tau(\alpha) =
\sigma(\tau(\alpha)/\alpha \cdot \alpha) =
\tau(\alpha)/\alpha\cdot\sigma(\alpha)$, karena hasil baginya terletak di
$K$), ke dalam sebuah grup siklik: jadi $\operatorname{Gal}$ siklik.
Kebalikannya adalah teori Kummer, yang tak akan kita perlukan (lihat
catatan di bawah). Untuk $K(\zeta_n)$: ia memecah \omterm{def:b3:galois:separable}{polinomial separabel} $X^n -
1$, dan $\sigma \mapsto a(\sigma)$ dengan $\sigma(\zeta_n) =
\zeta_n^{a(\sigma)}$ membenamkan grupnya ke dalam
$(\Z/n\Z)^\times$ yang abelian seperti pada \cref{thm:b3:galois:cyclotomicirred}
(keinjektifannya hanya menuntut $\zeta_n$ membangun akar
satuan yang terlibat).
\end{proof}

\begin{theorem}[Galois]\label{thm:b3:galois:solvable}
Misalkan $K$ berkarakteristik $0$ dan $P \in K[X]$ berlapangan
pemecah $L$. Jika $P$ \omterm{def:b3:galois:radical}{terselesaikan dengan radikal}, maka
$\operatorname{Gal}(L/K)$ merupakan \omterm{def:b3:groups:derived}{grup terselesaikan}. (Kebalikannya
juga benar; kita tak akan memerlukannya.)\index{teorema keterselesaian
Galois}
\end{theorem}

\begin{proof}
Langkah 1: \emph{perbesar menara \omterm{def:b3:galois:radical}{radikalnya} menjadi menara \omterm{def:b3:galois:galois}{Galois}.} Misalkan $L
\subseteq M$ dengan $M/K$ \omterm{def:b3:galois:radical}{radikal}, dengan eksponen akar $n_1, \dots,
n_r$ dan $n = n_1\cdots n_r$. Pertama tambahkan $\zeta_n$: menara
$K \subseteq K(\zeta_n) \subseteq M(\zeta_n)$ tetap \omterm{def:b3:galois:radical}{radikal}
(sebab $\zeta_n$ akar satuan: sebuah langkah \omterm{def:b3:galois:radical}{radikal}, $\zeta_n^n = 1$),
dan langkahnya setelah yang pertama terjadi atas lapangan yang memuat
akar satuan yang diperlukan. Berikutnya, gantilah $M(\zeta_n)$ dengan
komposit $N$ dari semua $\sigma(M(\zeta_n))$, dengan $\sigma$ menjelajahi
pembenaman-$K$ dari $M(\zeta_n)$ (yang berhingga banyaknya) ke dalam sebuah
\omterm{def:b3:galois:closure}{penutup aljabar} tetap: maka $N$ merupakan \omterm{thm:b3:galois:splitting}{lapangan pemecah} hasil kali
polinomial minimal sebuah himpunan pembangun (pada karakteristik $0$:
berhingga dan \omterm{def:b3:galois:separable}{separabel}), sehingga $N/K$ bersifat \omterm{def:b3:galois:galois}{Galois}; dan $N$ bersifat
\omterm{def:b3:galois:radical}{radikal} atas $K$: sebab setiap $\sigma(M(\zeta_n))$ \omterm{def:b3:galois:radical}{radikal} atas $K$
(terapkan $\sigma$ pada sebuah menara \omterm{def:b3:galois:radical}{radikal}), dan komposit perluasan
\omterm{def:b3:galois:radical}{radikal} juga \omterm{def:b3:galois:radical}{radikal} (sambungkan menaranya: jika $F'/K$ bersifat
\omterm{def:b3:galois:radical}{radikal} dengan menara yang menambahkan $\beta_j$, maka langkah bertipe
$F''(\beta_j)$ tetap \omterm{def:b3:galois:radical}{radikal} atas basis yang lebih besar mana pun).

Langkah 2: \emph{bacalah \omterm{def:b3:groups:derived}{keterselesaiannya} pada menara \omterm{def:b3:galois:galois}{Galois} itu.} Jadi andaikan
$L \subseteq N$, dengan $N/K$ \omterm{def:b3:galois:galois}{Galois} dan \omterm{def:b3:galois:radical}{radikal} bermenara $K
\subseteq K(\zeta_n) = E_0 \subseteq E_1 \subseteq \dots
\subseteq E_s = N$, dengan setiap $E_{i+1} = E_i(\sqrt[n_i]{a_i})$ dan
$\zeta_{n_i} \in E_0 \subseteq E_i$. Misalkan $G =
\operatorname{Gal}(N/K)$ dan $G_i = \operatorname{Gal}(N/E_i)$:
maka terbentuk rantai menurun $G \supseteq G_0 \supseteq G_1 \supseteq
\dots \supseteq G_s = \{e\}$. Setiap $E_{i+1}/E_i$ bersifat \omterm{def:b3:galois:galois}{Galois} dengan
grup siklik (\cref{lem:b3:galois:cyclicsteps}), sehingga menurut
teorema fundamental yang diterapkan pada \omterm{def:b3:galois:galois}{perluasan Galois} $N/E_i$
(\cref{thm:b3:galois:fundamental}(3), dengan grup lingkungannya $G_i$):
$G_{i+1} \trianglelefteq G_i$ dengan $G_i/G_{i+1} \cong
\operatorname{Gal}(E_{i+1}/E_i)$ yang siklik. Demikian pula $E_0/K$ bersifat
\omterm{def:b3:galois:galois}{Galois} dengan grup abelian $G/G_0$
(\cref{lem:b3:galois:cyclicsteps}). Rantai itu menampilkan $G$ sebagai
\omterm{def:b3:groups:derived}{grup terselesaikan} (\cref{prop:b3:groups:derived}). Akhirnya
$\operatorname{Gal}(L/K)$ merupakan \emph{\omterm{thm:b3:groups:quotient}{grup kuosien}} dari $G$: sebab $L/K$ bersifat
\omterm{def:b3:galois:galois}{Galois} ($P$ \omterm{def:b3:galois:separable}{separabel} pada karakteristik $0$) dan pembatasan $G
\to \operatorname{Gal}(L/K)$ bersifat surjektif
(\cref{thm:b3:galois:fundamental}(3) dengan $H =
\operatorname{Gal}(N/L)$); dan \omterm{thm:b3:groups:quotient}{grup kuosien} dari \omterm{def:b3:groups:derived}{grup terselesaikan} juga
\omterm{def:b3:groups:derived}{terselesaikan}.
\end{proof}

\begin{corollary}[Ketakterselesaian persamaan berderajat lima]
\label{cor:b3:galois:quintic}
Ada polinomial berderajat $5$ atas $\Q$ yang tidak
\omterm{def:b3:galois:radical}{terselesaikan dengan radikal}: misalnya $X^5 - 4X + 2$, yang grup
Galoisnya $S_5$ (\cref{exo:b3:galois:11}), sebuah grup tak \omterm{def:b3:groups:derived}{terselesaikan}
(\cref{cor:b3:groups:snnotsolvable}). Jadi tak ada rumus umum dalam
\omterm{def:b3:galois:radical}{radikal} untuk \omterm{def:b3:galois:extension}{derajat} $\geq 5$.\index{persamaan berderajat lima,
ketakterselesaian}
\end{corollary}

\begin{remark}\label{rem:b3:galois:converse}
Kebalikan \cref{thm:b3:galois:solvable} --- bahwa \omterm{def:b3:galois:galois}{grup Galois}
yang \omterm{def:b3:groups:derived}{terselesaikan} mengakibatkan \omterm{def:b3:groups:derived}{keterselesaian} dengan \omterm{def:b3:galois:radical}{radikal} --- dibuktikan dengan
menuruni deret \omterm{def:b3:groups:derived}{komutatornya} dan menunjukkan bahwa setiap perluasan siklik
(dengan cukup akar satuan) bersifat \omterm{def:b3:galois:radical}{radikal}, lewat \emph{resolven
Lagrange}; hal itu menjelaskan mengapa \omterm{def:b3:galois:extension}{derajat} $2, 3, 4$ mempunyai rumus:
sebab $S_2, S_3, S_4$ \omterm{def:b3:groups:derived}{terselesaikan}
(\cref{ex:b3:groups:solvableexamples}). Kita terima tanpa bukti pada
tingkat ini; pembahasan penuhnya termasuk kuliah magister, tetapi
\cref{exo:b3:galois:8} membuatnya konkret untuk kasus kubik.
\end{remark}

\section{Latihan}

\begin{exercise}[$\star$]\label{exo:b3:galois:1}
Tunjukkan $[\Q(\sqrt2, \sqrt3):\Q] = 4$, bahwa $\Q(\sqrt2 + \sqrt3) =
\Q(\sqrt2, \sqrt3)$, lalu hitunglah polinomial minimal
$\sqrt2 + \sqrt3$ atas $\Q$.
\end{exercise}

\begin{exercise}[$\star$]\label{exo:b3:galois:2}
Misalkan $\alpha = \sqrt[3]2$. Tunjukkan bahwa $\Q(\alpha)/\Q$ tidak
normal (tampilkan pembenaman $\Q(\alpha) \to \C$ yang petanya
bukan $\Q(\alpha)$), tentukan \omterm{thm:b3:galois:splitting}{lapangan pemecah} $L$ dari $X^3 -
2$ beserta $[L:\Q]$, lalu periksalah $\operatorname{Aut}_\Q(\Q(\alpha)) =
\{\mathrm{id}\}$: jadi pada perluasan yang bukan \omterm{def:b3:galois:galois}{Galois}, grup
automorfismanya bisa jauh lebih kecil daripada \omterm{def:b3:galois:extension}{derajatnya}.
\end{exercise}

\begin{exercise}[$\star$]\label{exo:b3:galois:3}
Bangunlah $\mathbb F_9$ sebagai $\mathbb F_3[X]/(X^2+1)$ lalu carilah
sebuah pembangun $\mathbb F_9^\times$. Daftarkan polinomial monik \omterm{def:b3:rings:divisibility}{tak tereduksi}
berderajat $1, 2, 3$ atas $\mathbb F_2$, lalu periksalah
$X^8 - X = X(X+1)(X^3+X+1)(X^3+X^2+1)$ atas $\mathbb F_2$.
\end{exercise}

\begin{exercise}[$\star\star$]\label{exo:b3:galois:4}
(a) Carilah semua akar primitif modulo $7$ dan modulo $11$ (yakni
pembangun $\mathbb F_7^\times$, $\mathbb F_{11}^\times$).
(b) Tunjukkan bahwa untuk $p$ ganjil, $x \in \mathbb F_p^\times$ merupakan
kuadrat jika dan hanya jika $x^{(p-1)/2} = 1$ (kriteria Euler), lalu peroleh kembali
kriteria bagi $-1$ pada \cref{pb:b3:rings:1}.
\end{exercise}

\begin{exercise}[$\star\star$]\label{exo:b3:galois:5}
Tunjukkan bahwa $\mathbb F_{p^m} \cap \mathbb F_{p^n} = \mathbb
F_{p^{\gcd(m,n)}}$ dan $\mathbb F_{p^m}\mathbb F_{p^n} = \mathbb
F_{p^{\operatorname{lcm}(m,n)}}$ di dalam sebuah \omterm{def:b3:galois:closure}{penutup aljabar}
tetap $\bar{\mathbb F}_p$, lalu uraikan
$\operatorname{Gal}(\mathbb F_{p^n}/\mathbb F_{p^m})$ untuk $m
\mid n$.
\end{exercise}

\begin{exercise}[$\star\star$]\label{exo:b3:galois:6}
Misalkan $I_d(q)$ banyaknya polinomial monik \omterm{def:b3:rings:divisibility}{tak tereduksi}
berderajat $d$ atas $\mathbb F_q$. Buktikan
\[
X^{q^n} - X \;=\; \prod_{d \mid n}\ \prod_{P \text{ tak tereduksi
monik, } \deg P = d} P ,
\qquad\text{sehingga}\qquad
q^n = \sum_{d \mid n} d\, I_d(q).
\]
Simpulkan $I_1, I_2, I_3, I_4$ secara eksplisit, dan $I_d(q) \geq 1$ untuk
setiap $d$ (sehingga perluasan $\mathbb F_{q^d}/\mathbb F_q$ ada sebagai
kuosien $\mathbb F_q[X]/(P)$ untuk setiap $d$).
\end{exercise}

\begin{exercise}[$\star\star$]\label{exo:b3:galois:7}
Tentukan $\operatorname{Gal}(\Q(\sqrt2,\sqrt3)/\Q)$ beserta
kisi lengkap lapangan antaranya. Pertanyaan yang sama untuk
\omterm{thm:b3:galois:splitting}{lapangan pemecah} $(X^2-2)(X^2-3)(X^2-6)$ --- apa yang Anda
perhatikan?
\end{exercise}

\begin{exercise}[$\star\star$]\label{exo:b3:galois:8}
(Persamaan kubik, diselesaikan oleh grupnya) Misalkan $P = X^3 + pX + q \in
\Q[X]$ \omterm{def:b3:rings:divisibility}{tak tereduksi} dengan akar $x_1, x_2, x_3$ dan \omterm{thm:b3:galois:splitting}{lapangan
pemecah} $L$. Misalkan $\delta = (x_1 - x_2)(x_1 - x_3)(x_2 - x_3)$ dan
$\Delta = \delta^2 = -4p^3 - 27q^2$ (terimalah identitas klasik ini
atau periksalah dengan menjabarkan fungsi simetrinya).
(a) Tunjukkan $\operatorname{Gal}(L/\Q) \cong A_3$ atau $S_3$, menurut
apakah $\Delta$ berupa kuadrat di $\Q$ atau bukan.
(b) Dengan $j = \zeta_3$, definisikan resolven Lagrange $u = x_1
+ jx_2 + j^2x_3$ dan $v = x_1 + j^2x_2 + jx_3$. Tunjukkan $u^3 + v^3
= -27q$ dan $uv = -3p$, lalu selesaikan untuk $u^3, v^3$: rumus Cardano
akan jatuh dengan sendirinya. Di manakah \omterm{def:b3:groups:derived}{keterselesaian} $S_3$ terpakai?
\end{exercise}

\begin{exercise}[$\star\star$]\label{exo:b3:galois:9}
Di $\Q(\zeta_5)$: tunjukkan bahwa satu-satunya sublapangan kuadratnya adalah
$\Q(\sqrt5)$, lewat jumlah Gauss $\eta_0 = \zeta_5 + \zeta_5^4$,
$\eta_1 = \zeta_5^2 + \zeta_5^3$: hitunglah $\eta_0 + \eta_1$ dan
$\eta_0\eta_1$, lalu simpulkan $\cos\frac{2\pi}5 = \frac{\sqrt5 -
1}4$. Simpulkan bahwa segi lima beraturan \omterm{def:b3:galois:constructible}{dapat dilukis}.
\end{exercise}

\begin{exercise}[$\star\star\star$]\label{exo:b3:galois:10}
Misalkan $K = \mathbb F_p(S, T)$ (fungsi rasional dalam dua
variabel tak tentu) dan $L = K(S^{1/p}, T^{1/p})$.
(a) Tunjukkan $[L:K] = p^2$ dan bahwa $\alpha^p \in K$ untuk setiap
$\alpha \in L$.
(b) Simpulkan bahwa $L/K$ \emph{tidak} sederhana: tak ada unsur primitif
yang tersedia --- jadi ketakseparabelan berakibat fatal bagi
\cref{thm:b3:galois:primitive}.
\end{exercise}

\begin{exercise}[$\star\star\star$]\label{exo:b3:galois:11}
Misalkan $P = X^5 - 4X + 2$ dan $G$ grup Galoisnya atas $\Q$,
yang beraksi pada $5$ akarnya.
(a) Tunjukkan bahwa $P$ \omterm{def:b3:rings:divisibility}{tak tereduksi}, lalu simpulkan $5 \mid \abs G$; simpulkan
bahwa $G$ memuat siklus-$5$ (lewat Cauchy,
\cref{thm:b3:groups:cauchy}).
(b) Tunjukkan, dengan mempelajari variasi $x \mapsto x^5 - 4x +
2$, bahwa $P$ mempunyai tepat $3$ akar real; lalu simpulkan bahwa
konjugasi kompleks terbatas menjadi sebuah transposisi di $G$.
(c) Tunjukkan bahwa subgrup $S_5$ yang memuat sebuah transposisi dan
sebuah siklus-$5$ pastilah $S_5$ \emph{(konjugasikan transposisinya dengan
pangkat siklus itu)}. Simpulkan $G \cong S_5$ dan, bersama
\cref{thm:b3:galois:solvable}, bahwa $P$ tidak \omterm{def:b3:galois:radical}{terselesaikan dengan
radikal}.
\end{exercise}

\begin{exercise}[$\star\star\star$]\label{exo:b3:galois:12}
(Kuartik dihedral) Misalkan $\alpha = \sqrt[4]2$ dan $L =
\Q(\alpha, \iu)$, yaitu \omterm{thm:b3:galois:splitting}{lapangan pemecah} $X^4 - 2$ atas $\Q$.
(a) Tunjukkan $[L : \Q] = 8$ dan bahwa $G =
\operatorname{Gal}(L/\Q)$ dibangun oleh $\sigma\colon
\alpha \mapsto \iu\alpha,\ \iu \mapsto \iu$ dan oleh konjugasi
kompleks $\tau$, dengan $\sigma^4 = \tau^2 = e$ dan
$\tau\sigma\tau = \sigma^{-1}$: jadi $G \cong D_4$.
(b) Daftarkan kisi subgrup $D_4$ (sepuluh subgrup) lalu
padankan masing-masing dengan lapangan tetapnya; periksalah khususnya bahwa
$\Q(\sqrt2)$, $\Q(\iu)$, $\Q(\iu\sqrt2)$ merupakan ketiga
sublapangan kuadratnya, lalu letakkan $\Q(\alpha)$,
$\Q(\iu\alpha)$, $\Q(\sqrt2, \iu)$.
(c) Lapangan antara mana yang bersifat \omterm{def:b3:galois:galois}{Galois} atas $\Q$? Cocokkan
jawaban Anda dengan \omterm{def:b3:groups:normal}{subgrup normal} $D_4$, lalu
jelaskan mengapa $\Q(\alpha)/\Q$ gagal sedangkan $\Q(\sqrt2)/\Q$
berhasil.
\end{exercise}

\section{Soal: Gauss dan segi 17 beraturan}

\begin{problem}[{Soal akhir pekan --- keterlukisan
segi 17}]\label{pb:b3:galois:1}
Pada 30 Maret 1796, Gauss yang berusia sembilan belas tahun menunjukkan bahwa
segi-$17$ beraturan \omterm{def:b3:galois:constructible}{dapat dilukis} --- kemajuan pertama atas
pertanyaan itu sejak zaman kuno. Kita susun ulang perhitungannya dengan
perkakas bab ini. Tetapkan $\zeta = \eu^{2\iu\pi/17}$, $L =
\Q(\zeta)$, $G = \operatorname{Gal}(L/\Q)$.

\medskip
\noindent\textbf{Bagian I --- Grup dan filtrasinya.}
\begin{enumerate}
  \item Berilah alasan: $[L:\Q] = 16$, $G \cong (\Z/17\Z)^\times$,
        yang siklik berorde $16$. Periksalah bahwa $3$ merupakan pembangun
        $(\Z/17\Z)^\times$ \emph{(hitunglah pangkat $3$
        modulo $17$: $3, 9, 10, 13, 5, 15, 11, 16, \dots$)}.
  \item Misalkan $\sigma \in G$ dengan $\sigma(\zeta) = \zeta^3$, dan
        $H_k = \langle \sigma^{2^k}\rangle$ untuk $k = 0, \dots,
        4$. Tunjukkan bahwa $G = H_0 \supset H_1 \supset H_2 \supset
        H_3 \supset H_4 = \{e\}$ dengan setiap indeks $[H_k :
        H_{k+1}] = 2$, dan bahwa lapangan tetapnya $\Q = L_0
        \subset L_1 \subset L_2 \subset L_3 \subset L_4 = L$
        membentuk menara perluasan kuadrat.
  \item Simpulkan \emph{secara a priori}, dengan memakai
        \cref{thm:b3:galois:wantzel}, bahwa $\zeta$ --- dan karenanya
        segi-$17$ --- \omterm{def:b3:galois:constructible}{dapat dilukis}. Sisa soal ini
        membuat menaranya eksplisit.
\end{enumerate}

\medskip
\noindent\textbf{Bagian II --- Periode berpanjang 8.}
Definisikan \emph{periode Gauss}
\[
\eta_0 = \sum_{k \text{ genap}} \zeta^{3^k \bmod 17}
= \zeta^{1} + \zeta^{9} + \zeta^{13} + \zeta^{15} + \zeta^{16} +
\zeta^{8} + \zeta^{4} + \zeta^{2},
\qquad
\eta_1 = \sum_{k \text{ ganjil}} \zeta^{3^k \bmod 17}.
\]
\begin{enumerate}[resume]
  \item Tunjukkan bahwa $\eta_0, \eta_1$ ditetapkan oleh $H_1$ dan ditukar
        oleh $\sigma$; lalu simpulkan $\eta_0, \eta_1 \in L_1$ dan bahwa
        keduanya merupakan dua akar sebuah polinomial kuadrat atas $\Q$.
  \item Hitunglah $\eta_0 + \eta_1 = -1$. Tunjukkan $\eta_0\eta_1 = -4$
        \emph{(setiap hasil kali $\zeta^a\zeta^b$ berupa suatu $\zeta^c$
        dengan $c \neq 0$; cacahlah berapa kali setiap $c$ muncul, atau
        berargumenlah bahwa hasil kalinya bilangan bulat rasional yang ditetapkan
        $G$, sama dengan jumlah atas seluruh $64$ hasil kali, lalu pakailah
        kenyataan bahwa setiap sisa tak nol muncul sama seringnya)}.
  \item Simpulkan $\eta_0 = \frac{-1 + \sqrt{17}}2$, $\eta_1 =
        \frac{-1-\sqrt{17}}2$ (kenali mana yang mana
        secara numerik: $\eta_0 \approx 1.56$), dan $L_1 =
        \Q(\sqrt{17})$.
\end{enumerate}

\medskip
\noindent\textbf{Bagian III --- Periode berpanjang 4 dan 2.}
Definisikan
\[
\beta_0 = \zeta + \zeta^{13} + \zeta^{16} + \zeta^{4},\quad
\beta_1 = \zeta^3 + \zeta^5 + \zeta^{14} + \zeta^{12},\quad
\beta_2 = \zeta^9 + \zeta^{15} + \zeta^{8} + \zeta^{2},\quad
\beta_3 = \zeta^{10} + \zeta^{11} + \zeta^{7} + \zeta^{6}.
\]
\begin{enumerate}[resume]
  \item Tunjukkan $\beta_0 + \beta_2 = \eta_0$, $\beta_1 + \beta_3 =
        \eta_1$, dan bahwa $\beta_0, \beta_2$ ditetapkan oleh $H_2$
        serta ditukar oleh $\sigma^2$.
  \item Hitunglah $\beta_0\beta_2 = -1$ dan $\beta_1\beta_3 = -1$
        \emph{(jabarkan: keenam belas eksponen yang diperoleh mencakup
        $1, \dots, 16$ tepat satu kali)}.
  \item Simpulkan $\beta_0 = \frac{\eta_0 + \sqrt{\eta_0^2 + 4}}2$
        (periksalah tandanya secara numerik: $\beta_0 \approx 2.05$) beserta
        rumus serupa bagi $\beta_1$; sehingga $L_2 =
        \Q(\beta_0)$, yang kuadrat atas $L_1$.
  \item Misalkan $\gamma_0 = \zeta + \zeta^{16} =
        2\cos\frac{2\pi}{17}$ dan $\gamma_1 = \zeta^{13} +
        \zeta^4$. Tunjukkan $\gamma_0 + \gamma_1 = \beta_0$ dan
        $\gamma_0\gamma_1 = \beta_1$, sehingga $\gamma_0 =
        \frac{\beta_0 + \sqrt{\beta_0^2 - 4\beta_1}}2$.
  \item Susunlah rantai rumus yang menyatakan
        $\cos\frac{2\pi}{17}$ lewat akar kuadrat bersarang, lalu berilah
        pemeriksaan desimalnya ($\cos\frac{2\pi}{17} \approx 0.93247$).
\end{enumerate}

\medskip
\noindent\textbf{Bagian IV --- Penutup.}
\begin{enumerate}[resume]
  \item Di manakah tepatnya argumen itu memakai bahwa $17$ merupakan
        \emph{prima Fermat} ($17 = 2^{2^2} + 1$)? Tunjukkan bahwa untuk
        bilangan prima $p$, segi-$p$ beraturan \omterm{def:b3:galois:constructible}{dapat dilukis} jika dan hanya jika $p
        = 2^{2^t} + 1$ untuk suatu $t$ \emph{(jika $p - 1 = 2^m$,
        tunjukkan bahwa $m$ sendiri haruslah pangkat $2$)}.
  \item Simpulkan daftar lengkap segi-$n$ beraturan yang \omterm{def:b3:galois:constructible}{dapat
        dilukis} untuk $n \leq 20$, dengan memakai kriteria Gauss--Wantzel
        pada \cref{cor:b3:galois:impossible}.
\end{enumerate}

\medskip
\noindent\textbf{Bagian V --- Jumlah Gauss dan resiprositas
kuadrat.} Periode pada Bagian II menyimpan sebuah harta. Untuk
bilangan prima ganjil $p$, \emph{simbol Legendre} $\bigl(\frac
ap\bigr)$ bernilai $+1$ jika $a$ kuadrat tak nol modulo $p$, $-1$ jika
bukan, dan $0$ jika $p \mid a$; \cref{exo:b3:galois:4}(b)
(kriteria Euler) memberikan $\bigl(\frac ap\bigr) \equiv
a^{(p-1)/2} \pmod p$, dan dari sini kemultiplikatifannya. Tulis $\zeta =
\eu^{2\iu\pi/p}$, $p^* = (-1)^{(p-1)/2}p$, lalu definisikan
\emph{jumlah Gauss}
\[
g \;=\; \sum_{a=1}^{p-1}\Bigl(\frac ap\Bigr)\zeta^a .
\]
\begin{enumerate}[resume]
  \item Tunjukkan $\sum_{a=1}^{p-1}\bigl(\frac ap\bigr) = 0$
        (sebanyak kuadratnya, sebanyak itu pula yang bukan kuadrat), lalu buktikan
        bentuk alternatifnya $g = \sum_{a=0}^{p-1}\zeta^{a^2}$
        \emph{(setiap kuadrat tak nol tercapai dua kali, dan
        $\sum_{a}\zeta^a = 0$)}. Untuk $p = 17$: kaitkan $g$ dengan
        periode pada Bagian II --- tunjukkan $g = \eta_0 - \eta_1$
        \emph{(sebab kuadrat modulo $17$ tepat berupa pangkat
        genap dari pembangun $3$)}.
  \item Buktikan $g^2 = p^*$: jabarkan
        \[
        g^2 = \sum_{a,b\neq0}\Bigl(\frac{ab}p\Bigr)
        \zeta^{a+b}
        = \sum_{c}\ \sum_{a \neq 0}\Bigl(\frac{a(c -
        a)}p\Bigr)\zeta^{c}
        \]
        (dengan menetapkan $b = c - a$), substitusikan $c - a = at$ untuk menilai
        jumlah di dalamnya sebagai $\bigl(\frac{-1}p\bigr)(p - 1)$ bila
        $c = 0$ dan $-\bigl(\frac{-1}p\bigr)$ bila tidak, lalu
        simpulkan dengan pertanyaan 14. Periksalah secara numerik: untuk $p =
        17$, $(\eta_0 - \eta_1)^2 = 17$ (Bagian II).
  \item Simpulkan $\sqrt{p^*} \in \Q(\zeta_p)$, lalu simpulkan
        bahwa sublapangan kuadrat \emph{tunggal} dari
        $\Q(\zeta_p)$ adalah $\Q(\sqrt{p^*})$ --- tunggal sebab
        $\operatorname{Gal}(\Q(\zeta_p)/\Q)$ bersifat siklik
        (\cref{thm:b3:galois:cyclotomicirred}) dan grup siklik
        mempunyai tepat satu subgrup berindeks $2$. (Setiap
        lapangan kuadrat terbenam di suatu lapangan siklotomik ---
        inilah kasus pertama teorema Kronecker--Weber,
        yang bentuk umumnya masih jauh di depan.)
  \item Sekarang misalkan $q \neq p$ bilangan prima ganjil lain. Dengan bekerja di
        dalam ring $\Z[\zeta]$ modulo $q$, buktikan
        \[
        g^q \equiv \Bigl(\frac qp\Bigr)\,g \pmod{q\Z[\zeta]}
        \]
        \emph{(mimpi mahasiswa baru: $(x + y)^q \equiv x^q + y^q$
        modulo $q$ di ring komutatif mana pun; lalu $g^q \equiv
        \sum_a\bigl(\frac ap\bigr)^q\zeta^{aq}$, indeks ulang $b =
        aq$ dan keluarkan $\bigl(\frac{q^{-1}}p\bigr) =
        \bigl(\frac qp\bigr)$)}.
  \item Di pihak lain, $g^q = g\,(g^2)^{(q-1)/2} =
        g\,(p^*)^{(q-1)/2}$; dengan memakai kriteria Euler modulo $q$,
        simpulkan $g^q \equiv \bigl(\frac{p^*}q\bigr)g
        \pmod{q\Z[\zeta]}$, lalu --- dengan mengalikan kedua
        ungkapan $g^q$ itu dengan $g$ dan memakai $g^2 = p^*$,
        yang terbalikkan modulo $q$ --- simpulkan
        \[
        \Bigl(\frac qp\Bigr) = \Bigl(\frac{p^*}q\Bigr) .
        \]
        \emph{(Mengapa kekongruenan antara bilangan bulat
        $\pm p^*$ modulo $q\Z[\zeta]$ mengakibatkan kesamaannya?
        Iriskan dengan $\Z$.)}
  \item Uraikan $\bigl(\frac{p^*}q\bigr) =
        \bigl(\frac{-1}q\bigr)^{(p-1)/2}\bigl(\frac pq\bigr)$
        dan $\bigl(\frac{-1}q\bigr) = (-1)^{(q-1)/2}$ untuk
        memperoleh \emph{hukum resiprositas kuadrat}:
        \[
        \Bigl(\frac pq\Bigr)\Bigl(\frac qp\Bigr)
        = (-1)^{\frac{p-1}2\cdot\frac{q-1}2} .
        \]
        Periksalah pada $(p, q) = (17, 3)$ dengan mendaftarkan kuadrat
        modulo $17$ dan modulo $3$, lalu pakailah untuk memutuskan dalam tiga
        baris apakah $x^2 \equiv 219 \pmod{383}$ \omterm{def:b3:groups:derived}{terselesaikan}
        ($383$ prima, $219 = 3\cdot73$).
\end{enumerate}

\medskip
\noindent\textbf{Bagian VI --- Mencacah polinomial \omterm{def:b3:rings:divisibility}{tak tereduksi}:
teorema bilangan prima untuk $\mathbb F_q[X]$.} Tetapkan sebuah pangkat prima
$q$ lalu misalkan $N_q(n)$ banyaknya polinomial \emph{monik \omterm{def:b3:rings:divisibility}{tak tereduksi}}
berderajat $n$ atas $\mathbb F_q$; ingat kembali dari
\cref{exo:b3:galois:6} faktorisasi $X^{q^n} - X$ dan
identitas $q^n = \sum_{d\mid n}d\,N_q(d)$, yang kini kita
balik, tafsirkan ulang, dan manfaatkan.
\begin{enumerate}[resume]
  \item (Kata) Sebutlah sebuah kata $w \in \mathbb F_q^n$
        \emph{primitif} jika ia bukan pangkat $u^{n/d} =
        u\cdots u$ dari kata $u$ yang lebih pendek secara sejati, lalu misalkan
        $A(d)$ banyaknya kata primitif berpanjang $d$.
        Tunjukkan bahwa setiap kata berpanjang $n$ secara tunggal merupakan pangkat
        sebuah kata primitif berpanjang $d \mid n$, sehingga
        $q^n = \sum_{d \mid n}A(d)$; lalu, dengan membandingkannya dengan
        \cref{exo:b3:galois:6}, simpulkan $A(d) = d\,N_q(d)$
        untuk setiap $d$, dan jelaskan kebetulan itu lewat sebuah
        bijeksi eksplisit: unsur $\alpha \in
        \mathbb F_{q^d}$ berderajat $d$ mempunyai \omterm{def:b3:groups:action}{orbit} \omterm{thm:b3:galois:finitefields}{Frobenius}
        $(\alpha, \alpha^q, \dots, \alpha^{q^{d-1}})$ yang berisi
        tepat $d$ unsur berbeda, dan unsur berderajat
        $d$ berkorespondensi $d$-ke-satu dengan polinomial \omterm{def:b3:rings:divisibility}{tak tereduksi} berderajat
        $d$.
  \item Buktikan \emph{rumus inversi Möbius}: jika $f(n) =
        \sum_{d\mid n}g(d)$ untuk setiap $n$, maka $g(n) =
        \sum_{d\mid n}\mu(d)\,f(n/d)$, dengan $\mu$ berupa
        fungsi Möbius ($\mu(m) = (-1)^{\#\text{faktor
        prima}}$ bila $m$ bebas kuadrat, dan $0$ bila tidak)
        \emph{(lema kuncinya: $\sum_{d \mid m}\mu(d) = 0$ untuk $m >
        1$ --- pasangkan pembagi yang memuat dan yang tidak memuat sebuah
        faktor prima tetap)}. Simpulkan
        \[
        N_q(n) = \frac1n\sum_{d \mid n}\mu(d)\,q^{n/d} .
        \]
  \item Tunjukkan $N_q(n) \geq \frac1n\bigl(q^n - 2q^{n/2}\bigr) > 0$
        untuk setiap $n \geq 1$: sebuah bukti baru bahwa $\mathbb F_{q^n}$
        ada untuk setiap $n$. Tafsirkan suku utamanya: polinomial
        monik acak berderajat $n$ bersifat \omterm{def:b3:rings:divisibility}{tak tereduksi} dengan
        peluang $\sim \frac1n$ --- padanan \omterm{prop:b3:galois:perfect}{sempurna} bagi
        teorema bilangan prima, dengan $\log x$ ditukar oleh $n$;
        periksalah secara numerik untuk $q = 2$, $n \leq 4$
        (\cref{exo:b3:galois:6} mendaftarkan cacahannya).
  \item Buktikan pendamping multiplikatif pertanyaan 21:
        \[
        \prod_{\substack{\pi \text{ monik tak tereduksi}\\ \deg\pi =
        n}}\pi \;=\; \prod_{d \mid n}
        \bigl(X^{q^d} - X\bigr)^{\mu(n/d)}
        \]
        \emph{(inversi Möbius di dalam grup abelian berisi
        fungsi rasional tak nol)}; periksalah dengan tangan untuk $q
        = 2$, $n = 2$: $(X^4 - X)/(X^2 - X) = X^2 + X + 1$.
\end{enumerate}

\medskip
\noindent\textbf{Bagian VII --- Dua penutup.}
\begin{enumerate}[resume]
  \item (Pelengkap kedua) Metode Bagian V juga menghitung
        $\bigl(\frac2q\bigr)$. Misalkan $\omega = \eu^{2\iu\pi/8}$
        dan $g = \omega + \omega^{-1}$. Tunjukkan $g^2 = 2$
        \emph{(sebab $\omega^2 = \iu$)}; lalu, untuk bilangan prima ganjil $q$,
        buktikan di $\Z[\omega]$ modulo $q$ bahwa
        \[
        g^q \equiv \omega^q + \omega^{-q} \pmod{q\Z[\omega]},
        \]
        dan bahwa ruas kanannya sama dengan $g$ bila $q \equiv \pm1
        \pmod 8$ dan $-g$ bila $q \equiv \pm3 \pmod 8$. Dengan membandingkannya
        dengan $g^q = g\,(g^2)^{(q-1)/2} \equiv
        \bigl(\frac2q\bigr)g$ seperti pada pertanyaan 18, simpulkan
        \[
        \Bigl(\frac2q\Bigr) = (-1)^{(q^2-1)/8},
        \]
        sambil memeriksa bahwa $(q^2 - 1)/8$ genap tepat ketika $q
        \equiv \pm1 \pmod 8$. Periksalah: $2$ merupakan kuadrat modulo $7$
        dan modulo $17$ (yakni $3^2$ dan $6^2$), tetapi bukan modulo $3$ maupun modulo
        $5$.
  \item (Fungsi zeta dari $\mathbb F_q[X]$) Buktikan
        identitas deret pangkat formal dalam $t$ berikut:
        \[
        \prod_{n \geq 1}\bigl(1 - t^n\bigr)^{-N_q(n)}
        = \frac1{1 - qt}
        \]
        \emph{(lewat ketunggalan faktorisasi atas semua polinomial monik \omterm{def:b3:rings:divisibility}{tak tereduksi}:
        jabarkan setiap faktornya sebagai deret geometri lalu cacahlah
        polinomial monik berderajat $n$)}. Perolehlah kembali identitas
        $q^m = \sum_{d \mid m}d\,N_q(d)$ pada
        \cref{exo:b3:galois:6} dengan mengambil logaritmanya. Periksalah
        koefisien $t^2$ dengan tangan untuk $q = 2$, lalu pakailah
        rumus pertanyaan 21 untuk menghitung $N_2(6) = 9$,
        sambil memeriksa $2^6 = 1\cdot2 + 2\cdot1 + 3\cdot2 +
        6\cdot9$.

\end{enumerate}
\end{problem}
